\chapter*{Abstract}

GNU/Linux has become the operating system of choice for many users, especially in the software development and system management fields. However, its increasing complexity presents significant challenges \cite{gonzalez2005analyzing}. Traditionally designed for terminal-based interactions, Linux requires innovative solutions to bridge the gap between its powerful capabilities and accessibility for a broader audience. This research explores the feasibility of employing a Large Language Model \cite{sheng2023s-loraLLM, devlin2019bertLLM}, specifically \GenModel, to create a cost-effective Autonomous Agent that enhances the functionalities of a Virtual Assistant \cite{liu2023whenLLM, marukatat2021textLLM, hu2021loraLLM, sheng2023s-loraLLM, devlin2019bertLLM} tailored for GNU/Linux users. The goal is to improve user experience by providing intelligent, context-aware recommendations and assisting with software management.

Beyond usability concerns, GNU/Linux also faces structural challenges such as distribution fragmentation and hardware compatibility \cite{zambrano2022implementacion}, which can hinder adoption and performance consistency across different environments. Despite these challenges, its open-source nature and extensive toolset make it a preferred platform for AI-driven applications. The integration of Artificial Intelligence (AI) into GNU/Linux has enabled advancements in automation, predictive analytics, and system optimization, reinforcing its role in cutting-edge technological development. As automation continues to shape software management, leveraging AI to streamline system interactions becomes increasingly necessary.

In this context, this research investigates the potential of \GenModel\ to develop a Virtual Assistant specifically designed for GNU/Linux users firstly apt-based distro. By incorporating AI-driven adaptability, the assistant aims to bridge the gap between complex system operations and user-friendly interactions. Unlike conventional assistants, this model will not only address general queries, but also provide system-specific guidance tailored to the user's installed software, ultimately enhancing accessibility, efficiency, and overall user satisfaction.


\textbf{Keywords: GNU/Linux, Large language models (LLM), Autonomous Agent, Virtual Assistant}